\section{Dobre, stare} % lang-pl

Eves: sekcja 7.2, 7.3. % lang-pl
Limit triangles. % lang-pl
Istnieje naturalna miara długości odcinków. % lang-pl
Czworokąt Saccheriego. % lang-pl
Suma kątów trójkąta jest $< \pi$. % lang-pl
Pole trójkąta, defekty. % lang-pl
Punkty idealne, ultraidealne, linie hiperrównoległe. % lang-pl

\begin{proposition}
[aksjomat Arystotelesa] % lang-pl
	Aksjomat Arystotelesa. % lang-pl
\index{aksjomat!Arystotelesa}%
\end{proposition}

\begin{proposition}
[lemat Proklusa] % lang-pl
	Prosta nie może przecinać tylko jednej z dwóch prostych równoległych. % lang-pl
\index{lemat!Proklusa}%
\end{proposition}

\begin{proposition}
[aksjomat Claviusa] % lang-pl
	Aksjomat Claviusa. % lang-pl
\index{aksjomat!Claviusa}%
\end{proposition}

\begin{proposition}
[aksjomat Clairauta] % lang-pl
	Aksjomat Clairauta. % lang-pl
\index{aksjomat!Clairauta}%
\end{proposition}

\begin{proposition}
[aksjomat Simsona] % lang-pl
	Aksjomat Simsona. % lang-pl
\index{aksjomat!Simsona}%
\end{proposition}

\begin{proposition}
[aksjomat Playfaira] % lang-pl
	Aksjomat Playfaira. % lang-pl
\index{aksjomat!Playfaira}%
\end{proposition}

Treść aksjomatu Playfaira poda szkocki matematyk John Playfair w podręczniku \emph{Elements of Geometry} z 1795 roku. % lang-pl
% % https://en.wikipedia.org/wiki/Playfair%27s_axiom
\index[persons]{Playfair, John}%

\begin{proposition}
[aksjomat Wallisa] % lang-pl
	Aksjomat Wallisa. % lang-pl
\index{aksjomat!Wallisa}%
\end{proposition}

\begin{proposition}
[aksjomat Bolyi] % lang-pl
	Aksjomat Bolyi. % lang-pl
\index{aksjomat!Bolyi}%
\end{proposition}

\begin{definition}
[czworokąt Saccheriego] % lang-pl
	Czworokąt Saccheriego. % lang-pl
% TODO: https://en.wikipedia.org/wiki/Lambert_quadrilateral
\index{czworokąt!Saccheriego}
\end{definition}

% TODO: https://www.deltami.edu.pl/2016/09/inne-swiaty-inne-geometrie/

Giovanni Girolamo Saccheri będzie uczniem Tommasa Cevy, brata Giovanniego, autora twierdzenia Cevy (w tej książce z numerem \ref{twierdzenie_cevy}). % lang-pl

O geometrii hiperbolicznej wspomina się w $\Delta_{24}^7$. % lang-pl

\begin{proposition}
	Jeśli geometria euklidesowa (na płaszczyźnie) jest niesprzeczna, to geometria Łobaczewskiego (na płaszczyźnie) też jest niesprzeczna. % lang-pl
\end{proposition}

Eves \cite[s. 347-353]{eves1_1972}. % lang-pl

W 1733 roku Giovanni Girolamo Saccheri spróbuje zbadać, co by było gdyby piąty postulat Euklidesa okazał się fałszywy. % lang-pl
Jego prace zapoczątkują systematyczne rozważania nad alternatywami dla geometrii euklidesowej. % lang-pl
W 1829 roku János Bolyai, Carl Friedrich Gauss oraz (niezależnie od nich) Nikołaj Łobaczewski stworzą geometrię hiperboliczną, czyli pierwszą w pełni rozwiniętą geometrię nieeuklidesową, w której piąty postulat Euklidesa nie obowiązuje. % lang-pl

Kolejnym osiągnięciem będzie model Beltramiego-Kleina dla geometrii hiperbolicznej, w którym pracujemy z punktami wewnątrz otwartego koła. % lang-pl
Proste w tym modelu to cięciwy łączące punkty z brzegu koła. % lang-pl
Wspomni o nim Eugenio Beltrami w 1868 roku, najpierw dla $n = 2$, a potem dla dowolnej liczby wymiarów. % lang-pl
% memoirs
Wcześniej, bo w 1859 roku, Arthur Cayley użyje dwustosunkowej definicji kąta (od Laguerre'a) by pokazać, że geometria euklidesowa może być zdefiniowana przy pomocy geometrii rzutowej. % lang-pl
Cayley zapozna się z pracami Beltramiego i Cayleya; w 1870 roku da wykład na seminarium u Weierstrassa, który zakończy pytaniem: czy między pomysłami Cayleya oraz Łobaczewskiegi istnieje powiązanie? % lang-pl

% na podstawie https://en.wikipedia.org/wiki/Beltrami–Klein_model
% gdzie pada
% and these essays proved the equiconsistency of hyperbolic geometry with ordinary Euclidean geometry

% model Kleina - Delta 2012-06

% Sferyczna osobno?
% https://en.wikipedia.org/wiki/Lexell's_theorem
\begin{proposition}
[twierdzenie Lexella] % lang-pl
	Wierzchołki trójkątów sferycznych o ustalonej podstawie $AB$ i polu powierzchni leżą na okręgu małym (powstałym przez przecięcie sfery płaszczyzną) przechodzącym przez punkty antypodalne do wierzchołków $A$, $B$ z podstawy. % lang-pl
\end{proposition}
% https://en.wikipedia.org/wiki/Spherical_trigonometry => Girard

Anders Johan Lexell napisze około 1777 roku pracę na temat tego twierdzenia, razem z dwoma dowodami: trygonometrycznym i geometrycznym. % lang-pl
Inni podadzą inne dowody: % lang-pl
Leonhard Euler w~1778, % lang-pl
Adrien-Marie Legendre w~1800, % lang-pl
Jakob Steiner w~1827, % lang-pl
Carl Friedrich Gauss w~1841, % lang-pl
Paul Serret w~1855, % lang-pl
Joseph-Émile Barbier w~1864. % lang-pl

% https://www.deltami.edu.pl/2012/06/w-rozumowaniach-byl-blad/
% https://www.deltami.edu.pl/2018/08/geometria-bolyaia-lobaczewskiego-co-to-jest-i-jak-ja-poznawac/

%
